% ============================================================
% KAPITEL 3 — ANFORDERUNGEN
% ============================================================

\chapter{Anforderungen}
\label{ch:anforderungen}

Die Anforderungen an das System stammen nicht aus einem vorgegebenen Lastenheft, sondern ergeben sich aus dem wissenschaftlichen Ziel der Arbeit und den Bedingungen, unter denen sie entstanden ist.
In diesem Kapitel werden zunächst die Quellen der Anforderungen benannt.
Danach folgen die funktionalen und nicht-funktionalen Anforderungen sowie die Zielkonflikte zwischen ihnen, und zuletzt wird abgegrenzt, was bewusst nicht gefordert ist.


% ------------------------------------------------------------
\section{Herleitung der Anforderungen}
\label{sec:anf-herleitung}

Da die Arbeit kein Auftragsprojekt mit externem Auftraggeber ist, gibt es keinen klassischen Anforderungskatalog.
Auch mit der Betreuung wurden keine formalen Anforderungen im Einzelnen festgelegt.
Die Anforderungen werden deshalb aus vier Quellen abgeleitet.

Das System soll in erster Linie den Einfluss der Repräsentationsform auf die natürlichsprachliche Szenenmanipulation messbar machen.
Aus den Forschungsfragen folgt deshalb, dass es mehrere Repräsentationsformen bereitstellen muss, die sich unter fairen Bedingungen vergleichen lassen.
Die immersive Zielumgebung, in der Datenschutz, geringe Latenz und die Unabhängigkeit von einer Netzverbindung zählen, spricht für lokale Modelle und verlangt, dass sich die Verarbeitungskette bis in die Brille übertragen lässt.

Damit die Ergebnisse belastbar sind, müssen die Messungen reproduzierbar und die Erfolgsbewertung objektiv sein.
Das betrifft vor allem die Kontrolle der Szenen und die Art, wie Erfolg gemessen wird.
Als feste Randbedingung kam die Hardware hinzu.
Für die Arbeit stand ein vom Institut bereitgestellter Arbeitsplatzrechner mit einer einzelnen Grafikkarte zur Verfügung, der die Größe der einsetzbaren Modelle begrenzt und damit mehrere der folgenden Anforderungen prägt.


% ------------------------------------------------------------
\section{Funktionale Anforderungen}
\label{sec:anf-funktional}

Zunächst muss das System einen natürlichsprachlichen Befehl entgegennehmen und in eine Veränderung der Szene umsetzen.
Die Operationen beschränken sich dabei auf das Erzeugen, Verschieben und Entfernen von Objekten.
Diese Beschränkung ist bewusst gewählt, da sich bereits mit diesen wenigen Operationen alle für die Untersuchung nötigen Arten von Referenzen abdecken lassen.

Die zentrale Anforderung ist, dass das System die Szene in den drei zu vergleichenden Repräsentationsformen bereitstellt, also strukturell, visuell und hybrid.
Damit der Vergleich dabei fair bleibt, müssen alle Varianten Objekte über dasselbe Interface ansprechen, d.\,h. über ein gemeinsames Kennungssystem.
Warum das für die Fairness so wichtig ist, wird in Abschnitt~\ref{sec:konz-fairness} zusammen mit den übrigen Entwurfsentscheidungen begründet.

Der Erfolg einer Manipulation muss sich objektiv bestimmen lassen.
Hierfür muss das System den tatsächlichen Zustand der Szene nach einer Manipulation auslesen können, sodass der Erfolg unabhängig von der Rückmeldung des Agenten ermittelt werden kann.
Das ist notwendig, da ein kleines Modell durchaus einen Erfolg melden kann, den es gar nicht erreicht hat.
Eine Messung anhand seiner eigenen Antwort wäre dann unzuverlässig.
Damit Messungen unter identischen Bedingungen wiederholt werden können, muss sich jedes Testszenario reproduzierbar aufbauen und zurücksetzen lassen.

Tabelle~\ref{tab:fa} fasst die funktionalen Anforderungen zusammen.
Die Spalte Herkunft gibt an, ob eine Anforderung aus den Forschungsfragen oder aus methodischen Gründen folgt.

\begin{table}[htbp]
  \centering
  \caption{Funktionale Anforderungen.}
  \label{tab:fa}
  \small
  \setlength{\tabcolsep}{5pt}
  \renewcommand{\arraystretch}{1.3}
  \begin{tabular}{@{}l p{9.4cm} l@{}}
    \toprule
    Nr. & Anforderung & Herkunft \\
    \midrule
    FA1 & Natürlichsprachliche Befehle in Szenenmanipulationen umsetzen, beschränkt auf Erzeugen, Verschieben und Entfernen von Objekten & RQ1, RQ2 \\
    FA2 & Die Szene in struktureller, visueller und hybrider Repräsentation bereitstellen & RQ1 \\
    FA3 & Objekte über ein gemeinsames Kennungssystem ansprechen, unabhängig von der Repräsentationsform & Validität \\
    FA4 & Den Erfolg einer Manipulation aus dem tatsächlichen Szenenzustand bestimmen & Validität \\
    FA5 & Testszenarien reproduzierbar aufbauen und zurücksetzen & Methodik \\
    \bottomrule
  \end{tabular}
\end{table}


% ------------------------------------------------------------
\section{Nicht-funktionale Anforderungen}
\label{sec:anf-nichtfunktional}

Die nicht-funktionalen Anforderungen betreffen keine einzelnen Funktionen, sondern Eigenschaften und Randbedingungen des Systems als Ganzes.

Aus der Zielumgebung folgt der lokale Betrieb der Modelle.
Da dafür nur die vorgegebene Hardware zur Verfügung steht, muss das System mit dem Grafikspeicher einer einzelnen Karte auskommen, und zwar so, dass Sprachmodell und visuelles Modell gleichzeitig im Speicher liegen.
Infrage kommen damit nur kleine Modelle.

Für die Aussagekraft der Untersuchung müssen sich unter identischer Konfiguration vergleichbare Ergebnisse einstellen, was kontrollierte und wiederholbar aufgebaute Szenen voraussetzt.
Da die eigentliche Zielumgebung die immersive Anwendung ist, muss dieselbe Verarbeitungskette außerdem sowohl am Desktop für die Messung als auch in der Brille laufen, ohne dass ihr Kern verändert wird.
Unterschiede lassen sich nur dann eindeutig auf die Repräsentation zurückführen, wenn sie die einzige Größe ist, die sich zwischen den Varianten ändert.
Tabelle~\ref{tab:nfa} fasst die nicht-funktionalen Anforderungen zusammen.

\begin{table}[htbp]
  \centering
  \caption{Nicht-funktionale Anforderungen.}
  \label{tab:nfa}
  \small
  \setlength{\tabcolsep}{5pt}
  \renewcommand{\arraystretch}{1.3}
  \begin{tabular}{@{}l p{9.4cm} l@{}}
    \toprule
    Nr. & Anforderung & Herkunft \\
    \midrule
    NFA1 & Modelle lokal betreiben & Zielumgebung \\
    NFA2 & Mit der vorgegebenen Hardware, einer einzelnen Grafikkarte, und Modellen geringer Größe auskommen & Hardware \\
    NFA3 & Messungen unter identischer Konfiguration reproduzierbar halten und die Konfiguration vollständig dokumentieren & Methodik \\
    NFA4 & Dieselbe Verarbeitungskette am Desktop und in der immersiven Umgebung betreiben & Zielumgebung \\
    NFA5 & Die Repräsentationsform als einzige zwischen den Varianten variierende Größe sicherstellen & Validität \\
    \bottomrule
  \end{tabular}
\end{table}

NFA5 ließ sich in der Umsetzung nur teilweise einhalten, was Abschnitt~\ref{sec:eval-validity} erläutert.


% ------------------------------------------------------------
\section{Zielkonflikte zwischen den Anforderungen}
\label{sec:anf-spannungen}

Die genannten Anforderungen lassen sich nicht alle gleichzeitig maximieren.
Am deutlichsten wird das bei Leistungsfähigkeit und Ressourcen.
Große Modelle würden natürlichsprachliche Befehle natürlich zuverlässiger umsetzen als die hier vorgegebenen kleinen lokalen Modelle.
Dieser Konflikt wird bewusst in Kauf genommen, da gerade das Verhalten unter dieser Beschränkung untersucht werden soll.

Ähnlich sieht es bei Fairness und Reproduzierbarkeit auf der einen und Praxisnähe auf der anderen Seite aus.
Das gemeinsame Kennungssystem und die einfachen, kontrolliert aufgebauten Szenen sind für einen sauberen und wiederholbaren Vergleich notwendig.
Gleichzeitig entfernen sie die Aufgabe ein Stück weit von einer realistischen Anwendung mit komplexen, realen Szenen.
In diesen Fällen hat die wissenschaftliche Kontrolle Vorrang, da das Ziel der Arbeit eine belastbare Erkenntnis und kein fertiges Produkt ist.
Die Übertragung in die reale Umgebung bleibt dem Demonstrator vorbehalten.


% ------------------------------------------------------------
\section{Abgrenzung des Funktionsumfangs}
\label{sec:anf-abgrenzung}

Passend zur Abgrenzung in Abschnitt~\ref{sec:ein-ziele} sind folgende Punkte bewusst nicht gefordert:
\begin{itemize}
  \item Produktreife: Es reicht, die untersuchten Repräsentationen und die Kette bis in die immersive Umgebung funktionsfähig umzusetzen, ohne Robustheit unter Dauerlast oder eine ausgereifte Bedienung.
  \item Weitere Operationen: Drehen oder Skalieren zur Laufzeit gehören nicht zur Untersuchung.
  \item Beliebige oder mehrdeutige Formulierungen: Die Befehle sind klar formuliert, Robustheit gegenüber umgangssprachlicher Vielfalt ist kein Ziel.
  \item Eine Nutzerstudie, da die Bewertung rein systembasiert aus dem Szenenzustand erfolgt.
  \item Übertragbarkeit auf andere Modelle oder reale, komplexe Szenen, da sich die Untersuchung auf ein Modellpaar und kontrolliert aufgebaute Szenen beschränkt.
\end{itemize}
